%% kapa: exam=1 round=22 year=2011 period=1 form=A subject=civil_law subject_ko=민법 seq=21 number=21
%% source: Q-Net 감정평가사 기출문제 (cst00309) artlSeq=5036790 file=extracted/1008/5036790/2011년 제22회 감정평가사 1차시험 기출문제/1교시.hwp
%% conversion: kordoc 4.24.51 -> tools/gen1.py (auto); review=proofread-sonnet; flags=none
\begin{question}{21}
민법상 법인에 관한 설명으로 옳은 것은?
\begin{choices}
  \item 임시이사는 대표권이 없다.
  \item 재단법인은 영리법인이 될 수 없다.
  \item 사단법인 사원의 지위는 양도할 수 없으나 상속할 수 있다.
  \item 사단법인 사원의 결의권은 평등하며, 정관으로 이를 달리 정할 수 없다.
  \item 법인 아닌 재단은 대표자 또는 관리인이 있는 경우에 소송상 당사자능력이 인정되지만, 등기능력은 인정되지 않는다.
\end{choices}
\end{question}
\answer{2}
